\documentclass{article}
\usepackage[utf8]{inputenc}
\usepackage{a4wide}

\begin{document}

\section*{เส้นระนาบ}

มีกราฟเส้นตรง (Line chart) ที่ประกอบด้วยจุดN{}จุด เชื่อมต่อกันด้วยส่วนของเส้นตรง

โจทย์กำหนดอาร์เรย์Yของเลขจำนวนเต็ม (ดัชนีของตัวแรกเท่ากับ1)

จุดที่Kมีตำแหน่งเป็น(K,Y{[}K{]}) ส่วนของเส้นตรงต่างๆ

ข้างต้นนั้นไม่มีเส้นแนวนอนนั่นคือไม่มีจุดสองจุดข้างกันที่มีค่าYเท่ากัน



หากเราลากเส้นตรงแนวนอนตัดกลุ่มของจุดเหล่านี้ให้หาจำนวนจุดสูงสุดที่เส้นตรงแนวนอนนี้สามารถตัดผ่านส่วนของเส้นตรงต่างๆ

ของNจุดนั้น



\hfill\break



{\def\LTcaptype{none} % do not increment counter

\begin{longtable}[]{@{}

  >{\raggedright\arraybackslash}p{(\linewidth - 4\tabcolsep) * \real{0.3333}}

  >{\raggedright\arraybackslash}p{(\linewidth - 4\tabcolsep) * \real{0.3333}}

  >{\raggedright\arraybackslash}p{(\linewidth - 4\tabcolsep) * \real{0.3333}}@{}}

\toprule\noalign{}

\begin{minipage}[b]{\linewidth}\centering

ข้อมูลเข้า

\end{minipage} & \begin{minipage}[b]{\linewidth}\centering

ข้อมูลออก

\end{minipage} & \begin{minipage}[b]{\linewidth}\centering

กราฟ

\end{minipage} \\

\midrule\noalign{}

\endhead

\bottomrule\noalign{}

\endlastfoot

\begin{minipage}[t]{\linewidth}\raggedright

7\\

1 2 1 2 1 3 2\\

\strut

\end{minipage} & 5 & \begin{minipage}[t]{\linewidth}\raggedright

\hfill\break

\includegraphics[width=6.82292in,height=2.52083in,alt={image.png}]{/public/upload/31a269c425.png}\\

\strut \\

\strut

\end{minipage} \\

\begin{minipage}[t]{\linewidth}\raggedright

5\\

2 1 3 4 5\\

\strut

\end{minipage} & 2 & \begin{minipage}[t]{\linewidth}\raggedright

\includegraphics[width=4.01042in,height=3.96875in,alt={image.png}]{/public/upload/abba35d9da.png}\\

\strut

\end{minipage} \\

\end{longtable}

}



\subsection*{Input}
\textbf{บรรทัดแรก}จำนวนเต็มหนึ่งค่าNกำหนดจำนวนจุดที่มี โดยที่N{∈}{[}2...500{]}



\textbf{บรรทัดที่สอง}จำนวนเต็มNค่า เป็นค่าY{[}K{]}โดยที่K\textbf{∈}{[}1...N{]}



\subsection*{Output}
จำนวนเต็มแสดงจำนวนจุดตัดสูงสุดที่เส้นตรงแนวนอนสามารถตัดผ่านส่วนของเส้นตรงต่างๆ ได้\\



\end{document}
